% ===================================================================
%  TABEL Nr. 2 — Het menselijk lichaam. — De pauze. — De spelen.
%  Traduction 2026 d'apres texto/io, controlee sur texto/fr et
%  texto/es. Consigne : voir texto/nl/10-tabel-01.tex.
% ===================================================================

%%K t02-apar-1 apar
\VUsaut{4.0mm}
\VUtitre{15.0pt}{60}{TABEL Nr. 2}
\VUsaut{2.0mm}
\VUtitre{11.4pt}{0}{Het menselijk lichaam. --- De pauze.}
\VUtitre{11.4pt}{0}{De spelen.}

%%K t02-tit-0 sub
\VUtitre{13.2pt}{0}{Het menselijk lichaam.}
\VUsaut{2.4mm}
\VUcentre{10.2pt}{0}{\textit{(Eenvoudige les in natuurwetenschap.)}}
\VUsaut{3.0mm}

%%K t02-tit-1 sub pos=t02-01-1
\VUpk{10.2pt}{40}{Het hoofd.}
\VUsaut{3.0mm}

%%K t02-01-1 p
1. --- De voornaamste delen van het menselijk lichaam zijn : het
\VUgras{hoofd}, de \VUgras{romp} en de \VUgras{ledematen}.

%%K t02-02-1 p
2. --- Het hoofd omvat twee delen : de \VUgras{schedel}, die de
\VUgras{hersenen} bevat en die het \VUgras{haar} bedekt, en het
\VUgras{gezicht}.

%%K t02-03-1 p
3. --- Op onze tekening zien wij noch de schedel noch de hersenen; maar wij
zien zeer duidelijk de belangrijke delen van het gezicht.

%%K t02-04-1 p
4. --- Ziehier eerst het \VUgras{voorhoofd}, dat een krachtig verstand aanduidt,
een altijd werkzaam denken. De \VUgras{slapen} zijn zeer vrij. Het
\VUgras{oog} is levendig en wijd geopend. Een dichte \VUgras{wenkbrauw} en lange
\VUgras{wimpers} beschutten het.

%%K t02-05-1 p
5. --- Ziehier bovendien de \VUgras{neus} en de \VUgras{neusgaten}. De neus
dient ons om te ruiken en om te ademen. Aan weerszijden van de neus zijn de
\VUgras{wangen}, verder weg de \VUgras{oren}. Het onderste deel van het gezicht
bestaat uit de \VUgras{mond} en de \VUgras{kin}.

%%K t02-06-1 p
6. --- Wij zien ze niet zo goed, want de \VUgras{snor} en de
\VUgras{bakkebaarden} verbergen ze voor ons. Maar wij weten, dat de mond de
twee \VUgras{lippen}, de \VUgras{bovenkaak} en de \VUgras{onderkaak} met hun
\VUgras{tanden}, en de \VUgras{tong} heeft. Met de mond spreken en eten wij. Met
de tong proeven wij het voedsel.

%%K t02-07-1 p
7. --- Het oog, het oor, de neus en de mond zijn, evenals de vingers, de
organen van de vijf zintuigen : het \VUgras{gezicht}, het \VUgras{gehoor}, de
\VUgras{reuk}, de \VUgras{smaak}, het \VUgras{gevoel}.

%%K t02-08-1 p
8. --- Het hoofd is dus een van de belangwekkendste delen van het menselijk
lichaam.

%%K t02-tit-2 sub
\VUsaut{4.4mm}
\VUpk{10.2pt}{40}{De romp.}
\VUsaut{3.0mm}

%%K t02-09-1 p
9. --- De \VUgras{hals} verbindt het hoofd met de romp. Hij omvat de
\VUgras{keel} en de \VUgras{nek}.

%%K t02-10-1 p
10. --- De romp omvat ook veel wezenlijke organen. In de \VUgras{borst} zijn de
\VUgras{longen}, waarmee wij ademen, en het \VUgras{hart} dat het middelpunt van
het leven is. Wanneer het hart ophoudt te kloppen, sterft het levende wezen.

%%K t02-11-1 p
11. --- De \VUgras{ribben} houden de borst overeind.

%%K t02-12-1 p
12. --- De overige delen van de romp zijn de \VUgras{zijde}, de \VUgras{rug}, de
\VUgras{schouders} en de \VUgras{buik}. Wij gaan op de zijde of op de rug liggen
om te slapen, wij dragen de lasten op onze schouder.

%%K t02-13-1 p
13. --- In de buik zijn de \VUgras{maag} en de \VUgras{darmen}. Zij dienen ons
om het voedsel te verteren.

%%K t02-tit-3 sub
\VUsaut{4.4mm}
\VUpk{10.2pt}{40}{De ledematen.}
\VUsaut{3.0mm}

%%K t02-14-1 p
14. --- Wij hebben vier \VUgras{ledematen} : twee \VUgras{armen} en twee
\VUgras{benen}. Met de armen nemen, houden, heffen en verzamelen wij de
voorwerpen; met de benen staan, lopen en rennen wij.

%%K t02-15-1 p
15. --- De \VUgras{schouder} verbindt de arm met het lichaam. Een zeer stevige
\VUgras{spier}, de biceps, beweegt de arm die de \VUgras{elleboog} met de
\VUgras{onderarm} verbindt. De \VUgras{pols} vormt het gewricht van de
\VUgras{hand}, die eindigt met de \VUgras{vingers} en de \VUgras{nagels}. Op de
hand onderscheiden wij een \VUgras{ader} (en een slagader) waarin het
\VUgras{bloed} stroomt.

%%K t02-16-1 p
16. --- De delen van het been zijn : het \VUgras{dijbeen}, de \VUgras{knie} en
de \VUgras{knieholte}, de \VUgras{kuit}, waarvan het \VUgras{vlees} door de
\VUgras{huid} bedekt is, de \VUgras{enkel} en de \VUgras{voet}. De
\VUgras{beenderen} van het been zijn zeer dik en zeer sterk. Aan het uiteinde
van de voet zijn de \VUgras{tenen}. De overige delen zijn de \VUgras{voetzool}
en de \VUgras{hiel}.

%%K t02-17-1 p
17. --- Zo luidt de bijna volledige beschrijving van het menselijk lichaam.

%%K t02-17-2 p
\textit{Nota.} --- \textit{Met behulp van de uitdrukking : « ik heb pijn aan...
(het hoofd enz.) » zal de onderwijzer deze hele woordenschat opnieuw kunnen
gebruiken uit het oogpunt van de goede of slechte} \VUgras{lichamelijke
toestand.}

%%K t02-tit-4 sub
\VUsaut{5.0mm}
\VUpk{10.2pt}{40}{Gymnastiek.}
\VUsaut{2.4mm}
\VUcentre{10.2pt}{0}{\textit{(Verhaal van een van de leerlingen.)}}
\VUsaut{3.0mm}

%%K t02-18-1 p
18. --- Om lenig, behendig en gezond te zijn, moeten wij onze benen en onze
armen oefenen. Daarom spelen wij en beoefenen wij de \VUgras{gymnastiek}.

%%K t02-19-1 p
19. --- Tijdens de week vinden deze twee oefeningen plaats op de binnenplaats
van het lyceum (zie tabel nr. 1). Maar op donderdag en op zondag gaan wij naar
een groot \VUgras{grasveld}, waar wij ruimte hebben om te rennen en ons te
vermaken.

%%K t02-20-1 p
20. --- Onze meesters en onze ouders vergezellen ons daar soms; maar een
\VUgras{leraar gymnastiek} \textsuperscript{(12)} leidt wel degelijk de
oefeningen.

%%K t02-21-1 p
21. --- Op het grasveld, links van de ingang, staan de
\VUgras{gymnastiektoestellen} \textsuperscript{(1)}, wij klimmen langs de
\VUgras{ladder} \textsuperscript{(2)}, wij hangen ondersteboven aan de
\VUgras{ringen} \textsuperscript{(3)} en aan de \VUgras{trapeze}
\textsuperscript{(4)}, wij hijsen ons met de handen op tot de top van de
\VUgras{touwladder} \textsuperscript{(5)} of van het \VUgras{knoopentouw}
\textsuperscript{(6)}.

%%K t02-22-1 p
22. --- Intussen springen onze kameraden over het \VUgras{houten paard}
\textsuperscript{(7)} of de \VUgras{brug met gelijke leggers}
\textsuperscript{(11)}. Anderen \VUgras{boksen} \textsuperscript{(8)} of
\VUgras{schermen} \textsuperscript{(9)} met een masker en met een
\VUgras{floret}. Ten slotte \VUgras{heffen} anderen \VUgras{halters}
\textsuperscript{(10)}.

%%K t02-tit-5 sub
\VUsaut{4.6mm}
\VUpk{10.2pt}{40}{De spelen.}
\VUsaut{3.0mm}

%%K t02-23-1 p
23. --- Rondom de turners spelen jongens en meisjes verschillende
\VUgras{spelen}. De meisjes vermaken zich met hun \VUgras{hoepel}
\textsuperscript{(14)}; zij springen aan het \VUgras{touw}
\textsuperscript{(13)} of nodigen hun broers en neven uit voor een partij
\VUgras{croquet} \textsuperscript{(15)}. Zij hebben er plezier in de
\VUgras{bal} van de tegenstander weg te stoten, juist op het ogenblik dat hij
denkt gemakkelijk door het poortje te gaan!

%%K t02-24-1 p
24. --- De jongens houden meer van de spelen die spierkracht vragen. Graag
spelen zij \VUgras{kegelen} \textsuperscript{(16)} waarbij zij de kegels
behendig omgooien met een \VUgras{bal} \textsuperscript{(17)}. Zij versmaden
het ook niet met een \VUgras{tol} \textsuperscript{(18)} en een
\VUgras{bilboquet} \textsuperscript{(19)} of zelfs met de \VUgras{kikker}
\textsuperscript{(20)} te spelen. Maar hun lievelingsspelen zijn wel de
\VUgras{knikkers} \textsuperscript{(21)} en het \VUgras{kaatsen tegen de muur}
\textsuperscript{(22)}, want de muur van de \VUgras{broeikas}
\textsuperscript{(26)} is zeer geschikt om de lichte bal te laten terugspringen.

%%K t02-25-1 p
25. --- De kleine Jan, die op zijn \VUgras{stelten} \textsuperscript{(24)}
loopt, is zeer behendig en weet zeer goed zijn evenwicht te bewaren. Naast hem
spelen enkele kameraden \VUgras{verstoppertje} \textsuperscript{(25)} in die
broeikas en achter de \VUgras{heg}. Anderen houden meer van het
\VUgras{slagraden} \textsuperscript{(28)} of van de \VUgras{pluimbal}
\textsuperscript{(29)} die van het ene \VUgras{racket} naar het andere vliegt.

%%K t02-noto-1 noto
\VUnotes{143.2mm}{%
(*) De kinder- of jongensspelen hebben dikwijls zuiver nationale namen. Wij
hebben gehandeld om ze in Ido zo goed mogelijk te benoemen, hetzij door ons te
laten inspireren door de handeling zelf, waardoor zij gekenmerkt worden, hetzij
door de nationale naam in dit of dat land te kiezen, wanneer die niet al te
onbillijk is. Bijv. : het \VUgras{tonnetjesspel} (in het Duits en in het Frans)
is het \VUgras{kikkerspel} \textsuperscript{(20)} (in het Engels) waarin de
spelers hun ronde schijfjes trachten te werpen door de mond van de metalen
kikker die met open bek op de doorboorde kist (het tonnetje) staat. Het
\VUgras{schoudersprongspel} \textsuperscript{(30)} verschilt van het
\VUgras{bokspringen} alleen hierin : om over het hoofd te springen steunen de
spelers de handen op de schouders van hun makker, terwijl zij bij het
« \VUgras{bokspringen} » de handen alleen op zijn rug steunen. --- Of ook zijn
enkele deelnemers gebukt terwijl de anderen over hun rug zullen springen met de
handen op de schouder gesteund; de eersten moeten de schok verdragen zonder door
te buigen, de tweeden moeten springen zonder het evenwicht te verliezen (Zie de
tabel zelf).%
}

%%K t02-26-1 p
26. --- Midden op het grasveld staan twee bomen. Aan de voet van een van beide
hebben zeven of acht jongens een partij \VUgras{schoudersprong}
\textsuperscript{(30)} (*) op touw gezet. Niet in alle landen kent men dat
spel; maar iedereen weet wat het \VUgras{bokspringen} is, dat de jongens deze
keer niet spelen.

%%K t02-tit-6 sub
\VUsaut{5.0mm}
\VUpk{10.2pt}{40}{De spelen in de open lucht.}
\VUsaut{3.4mm}

%%K t02-27-1 p
27. Het \VUgras{blindemannetje} \textsuperscript{(31)} is ook zeer vermakelijk;
meisjes en jongens doen om beurten de \VUgras{blinddoek} voor hun ogen en
grijpen de onhandigen die zich laten vangen.

%%K t02-28-1 p
28. --- Midden op het grasveld, ziehier een zeer Frans maar wat ruw spel : het
is de \VUgras{berin} \textsuperscript{(32)} die veel luidruchtiger is dan het zo
rustige spel van de \VUgras{vlieger} \textsuperscript{(33)} of zelfs dan het
Engelse spel \VUgras{tennis} \textsuperscript{(34)}.

%%K t02-29-1 p
29. --- Die laatste sport is de vreugde van de grote leerlingen. Onze leraren
zelf beoefenen hem graag. Hij vraagt grote behendigheid met lenigheid en hij
bevalt mij zeer. Ik laat aan de meisjes de spelen met de \VUgras{schommel}
\textsuperscript{(35)}.

%%K t02-30-1 p
30. --- Ik houd niet van een ander Engels spel dat misschien gevaarlijk zou
worden.

%%K t02-30-1b p
Ik bedoel het \VUgras{voetbal} \textsuperscript{(36)} dat mijn oudere broer
verheugt. Ik geef de voorkeur aan ons oude \VUgras{barrenspel}
\textsuperscript{(38)} dat even gezond en werkelijk minder ruw is.

%%K t02-tit-7 sub
\VUsaut{4.6mm}
\VUpk{10.2pt}{40}{Fietsen en balspel.}
\VUsaut{3.0mm}

%%K t02-31-1 p
31. --- Ook rijd ik graag met de \VUgras{fiets} \textsuperscript{(37)} om het
grasveld heen.

%%K t02-32-1 p
32. --- Volgend jaar zal ik het \VUgras{paardrijden} \textsuperscript{(39)}
leren. Zo mooi is een paard dat sierlijk stapt of draaft, en dat in galop over
de hindernissen springt!

%%K t02-33-1 p
33. --- In de zomer leren wij de \VUgras{zwemkunst} \textsuperscript{(40)}. Wij
duiken met genot en baden in het heldere en frisse water van de rivier. Soms
laten wij ten slotte een papieren \VUgras{ballon} \textsuperscript{(41)} op die
statig in de dampkring opstijgt, zodra hij met warme lucht gevuld is.

%%K t02-34-1 p
34. --- Er zijn ook nog veel andere spelen, maar ik zou er geen einde aan
vinden ze op te sommen en men beseft aan de hand van dit overzicht, dat men zich
op donderdag niet erg verveelt op ons mooie grasveld.

%%K t02-34-2 p
N.-B. --- \textit{De onderwijzer zal dat hoofdstuk gemakkelijk aanvullen; door
elk van de belangrijkste spelen uitvoeriger te beschrijven.}
\VUsaut{6.0mm}
\VUfilet{22mm}
